\section{Введение}

С ростом объёмов обрабатываемых данных в биоинформатике, астрономии, машинном обучении и других областях задачи эффективного хранения и доступа к информации становятся всё более актуальными. Классические алгоритмы сжатия (GZIP, BZIP2) требуют полной декомпрессии файла для чтения или модификации, что неприемлемо для многих приложений. В биоинформатике существуют специализированные решения, обеспечивающие произвольный доступ к сжатым данным (CRAM), однако они привязаны к конкретному типу данных (геном). Также существуют более универсальные инструменты (seekable zstd, zran), однако они поддерживают только чтение, без возможности эффективной перезаписи.

Однако в некоторых специализированных задачах так же возникает необходимость не только в операциях чтения и записи, но и операциях вставки и удаления из середины файла (с изменением размера файла). Существующие решения так же не предоставляют эффективной реализации подобных операций.

Помимо этого, выбор алгоритма сжатия существенно влияет на степень сжатия и скорость работы. Существуют как универсальные алгоритмы (Zlib, LZ4, Zstd, Brotli), так и специализированные, разработанные под конкретные типы данных (например, геномные или астрономические \textbf{TODO}). Хотелось бы иметь инструмент, которые не завязан на конкретном алгоритме сжатия, а предоставляющий возможность выбора любого алгоритма сжатия, в том числе и подключения своего.

Целью данной работы является разработка библиотеки \texttt{Compio}, обеспечивающей прозрачное сжатие с произвольным доступом на чтение, запись, вставку и удаление, с интерфейсом, аналогичным \texttt{stdio}, поддержкой нескольких файлов в одном архиве и возможностью подключения произвольных алгоритмов сжатия.

В работе проведён анализ существующих систем, предложена архитектура файла-контейнера с блочным сжатием и индексацией на основе B-дерева, реализовано кэширование узлов и данных, разработано API для языка C, выполнена экспериментальная оценка производительности и сравнение с некоторыми альтернативными решениями.